\documentclass[12pt]{amsart}
\usepackage{mathrsfs,amssymb, amsmath, mathdots, xcolor, tikz-cd}


\textwidth = 160mm
\textheight = 240mm
\evensidemargin=0mm
\oddsidemargin=0mm
\hoffset=-1mm
\voffset=-22mm
\parskip =1mm
\parindent = 5mm
\linespread{1.2}
\footskip10mm
\pagestyle{plain}

\newtheorem{theorem}{Theorem}[section]
\newtheorem{lemma}[theorem]{Lemma}
\newtheorem{corollary}[theorem]{Corollary}
\newtheorem{proposition}[theorem]{Proposition}


\numberwithin{equation}{section}

\theoremstyle{definition}
\newtheorem{definition}[theorem]{Definition}
\newtheorem{example}[theorem]{Example}
\newtheorem{construction}[theorem]{Construction}
\newtheorem{problem}[theorem]{Problem}
\newtheorem{remark}[theorem]{Remark}
\newtheorem{question}[theorem]{Question}
\newtheorem{conjecture}[theorem]{Conjecture}
\newtheorem{convention}[theorem]{Convention}
\newtheorem{notation}[theorem]{Notation}
\newtheorem{defprop}[theorem]{Definition-Proposition}

\DeclareSymbolFont{script}{U}{eus}{m}{n}
\DeclareMathSymbol{\Wedge}{0}{script}{"5E}

\newcommand{\laplacian}{\Delta}
\newcommand{\cA}{\mathscr{A}} 
\newcommand{\grad}{\operatorname{gradient \ term}} 
\newcommand{\cG}{\mathcal{G}} 
\newcommand{\cI}{\mathcal{I}} 
\newcommand{\RR}{\mathbb{R}} 
\renewcommand{\AA}{\mathbb{A}} 
\newcommand{\DD}{\mathbb{D}}
\newcommand{\NN}{\mathbb{N}} 
\newcommand{\hH}{\mathbb{H}} 
\newcommand{\CC}{\mathbb{C}} 
\newcommand{\cC}{\mathcal{C}} 
\newcommand{\cK}{\mathcal{K}} 
\newcommand{\ZZ}{\mathbb{Z}} 
\newcommand{\cX}{\mathfrak{X}} 
\newcommand{\cY}{\mathcal{Y}} 
\newcommand{\cV}{\mathcal{V}}
\newcommand{\GGm}{\mathbb{G}_m}
\newcommand{\GG}{\mathbb{G}}
\newcommand{\RRPP}{\mathbb{RP}} 
\newcommand{\PP}{\mathbb{P}} 
\newcommand{\fF}{\mathcal{F}} 
\newcommand{\cL}{\mathcal{L}} 
\newcommand{\gG}{\mathcal{G}} 
\renewcommand{\AA}{\mathbb{A}} 
\newcommand{\QQ}{\mathbb{Q}}
\newcommand{\cF}{\mathcal{F}} 
\newcommand{\cE}{\mathcal{E}} 
\newcommand{\cR}{\mathcal{R}}
\newcommand{\cU}{\mathcal{U}}  
\newcommand{\cP}{\mathcal{P}} 
\newcommand{\Frac}{\operatorname{Frac}}
\newcommand{\Proj}{\operatorname{Proj}}
\newcommand{\trdeg}{\operatorname{trdeg}}
\newcommand{\id}{\operatorname{id}}
\newcommand{\htt}{\operatorname{ht}}
\newcommand{\im}{\operatorname{im}}
\newcommand{\Bl}{\operatorname{Bl}}
\newcommand{\Gr}{\mathrm{Gr}}
\newcommand{\Aut}{\operatorname{Aut}}
\newcommand{\bsim}{\sim_{\operatorname{bir}}}
\newcommand{\bir}{\operatorname{bir}}
\newcommand{\dR}{\operatorname{dR}}
\newcommand{\dB}{\operatorname{dB}}
\newcommand{\HH}{\operatorname{H}}
\newcommand{\cH}{\mathcal{H}}
\newcommand{\bH}{\mathbb{H}}
\newcommand{\Exc}{\operatorname{Exc}}
\newcommand{\tr}{\operatorname{tr}}
\newcommand{\Ric}{\operatorname{Ric}}
\newcommand{\Pic}{\operatorname{Pic}}
\newcommand{\NS}{\operatorname{NS}}
\newcommand{\PGL}{\operatorname{PGL}}
\newcommand{\Tor}{\operatorname{Tor}}
\newcommand{\dvol}{d\mathrm{vol}}
\newcommand{\Sch}{\mathrm{Sch}}
\newcommand{\pr}{\mathrm{pr}}
\newcommand{\Hom}{\mathrm{Hom}}
\newcommand{\Hilb}{\mathrm{Hilb}}
\newcommand{\Schet}{\mathrm{Sch}_{\text{\'et}}}
\newcommand{\Schfpqc}{\mathrm{Sch}_{\text{fpqc}}}
\newcommand{\Spec}{\mathrm{Spec}}
\newcommand{\Supp}{\mathrm{Supp}}
\newcommand{\coker}{\mathrm{coker}}
\renewcommand{\div}{\operatorname{div}}
\let\bar\overline
\let\phi\varphi
\let\tilde\widetilde
\let\hat\widehat
\let\epsilon\varepsilon

\newcommand\mapsfrom{\mathrel{\reflectbox{\ensuremath{\mapsto}}}}


% you can add more abbreviations of your own here

\begin{document}

\title{Hilbert Scheme}

\author{Tyson Klingner}


\date{\today}  % the date on which the report was finished

\maketitle 

\section{Introduction}

To understand the geometry of a given algebraic variety $X$ we often need to understand some of its subschemes. 
For example, if $X = \PP^n$, then all of its smooth degree $d$ hypersurfaces are diffeomorphic. 
When studying subschemes one quickly sees that subschemes comes in families, e.g.,
\[
V(xy-tz^2) \subseteq \PP^2_{x,y,z} \times \AA^1_t
\]
defines a family of conics in $\PP^2.$
Hence, to understand subschemes we will study them in families. 
We achieve this via the {\em Hilbert functor}. 
Once we introduce the Hilbert functor we will show that the Hilbert functor is represetable by a projective scheme, which we call the {\em Hilbert scheme}.
In what proceeds we work over $k= \bar{k}$. 
Every construction extends to the relative setting, but we work over $k$ for simplicity.

\section{Hilbert Functor}

Let $X/k$ be a projective variety and let $\mathcal{O}_X(1)$ be a fixed very ample line bundle. Before defining the Hilbert functor we will introduce the {\em Hibert polynomial.}

\begin{defprop}
Let $\cF$ be a coherent sheaf on $X$, then there is a numerical polynomial $P(z) \in \QQ[z]$, such that $\chi(\cF(n)) = P_{\cF}(n)$ for every $n \in \ZZ.$
We call $P_{\cF}$ the {\em Hilbert polynomial} of $\cF$ with respect to $\mathcal{O}_X(1).$
If $Z \subset X$ is a closed subscheme, then we define the {\em Hilbert polynomial} of $Z$ to be $P_{\mathcal{O}_Z}.$ 
\end{defprop}

\begin{proposition}[Properties of Hilbert Polynomials]
    Here we collect some basics properties of Hilbert polynomials
\begin{enumerate}
    \item Suppose there is a short exact sequence of coherent sheaves
    \[
    0 \to \cE \to \cF \to \cG \to 0.
    \]
    Then, 
    \[
    P_{\cF}(n) = P_{\cE}(n) + P_{\cG}(n)
    \]
    \item 
\end{enumerate}
\end{proposition}    

\begin{example}
Suppose $X = \PP^m$. Then, $\chi(\mathcal{O}_{\PP^m}(n)) = {{m+n}\choose m}$ for $n > > 0$ so the Hilbert polynomial of $\PP^m$ is given by
\[
P_{\PP^m}(n) = {{m+n}\choose{m}}, \ \ \ n > > 0
\]
In particular,
\[
P_{\PP^2}(n) = \frac{(n+1)(n+2)}{2}, \ \ \ n > > 0.
\]
If $C \subseteq \PP^2$ is a degree $d$ curve, then there is a short exact sequence
\[
0 \to \mathcal{O}_{\PP^2}(n-d) \to \mathcal{O}_{\PP^2}(n) \to \mathcal{O}_C(n) \to 0
\]
and since $P_C(n) = P_{\PP^2}(n) - P_{\PP^2}(n-d)$, it follows that
\[
P_C(n) = dn - \frac{1}{2}d(d-3), \ \ \ n > > 0.
\]
\end{example}    



\end{document}